\documentclass[aspectratio=169]{beamer}
% 16:9 aspect ratio for modern projectors

% Theme and color scheme
\usetheme{Madrid}
\usecolortheme{default}

% Custom colors matching the design recommendations
\definecolor{primaryblue}{RGB}{0,51,102}      % Deep blue (mathematical/technical)
\definecolor{secondarygreen}{RGB}{0,128,0}   % Green (sports/football)
\definecolor{accentorange}{RGB}{255,140,0}    % Orange (highlights/key findings)
\definecolor{lightgray}{RGB}{245,245,245}     % Light grey background

\setbeamercolor{structure}{fg=primaryblue}
\setbeamercolor{title}{fg=primaryblue}
\setbeamercolor{frametitle}{fg=primaryblue}
\setbeamercolor{block title}{bg=primaryblue,fg=white}
\setbeamercolor{block body}{bg=lightgray,fg=black}

% Packages
\usepackage{amsmath}
\usepackage{amssymb}
\usepackage{graphicx}
\usepackage{booktabs}
\usepackage{multirow}
\usepackage{xcolor}
\usepackage{tikz}
\usepackage{fontawesome5}

% Title information
\title[Football-TDA]{Topological Data Analysis for Dynamic Team Behaviour}
\subtitle{Multi-Scale Persistent Homology in Competitive 22-Body Systems}
\author{GPS-TDA Research Team}
\institute{Mathematical Foundations for Sports Analytics}
\date{December 2024}

% Remove navigation symbols
\setbeamertemplate{navigation symbols}{}

% Custom commands
\newcommand{\highlight}[1]{\textcolor{accentorange}{\textbf{#1}}}
\newcommand{\scaleitem}[2]{\textbf{#1}: #2}

\begin{document}

% Title slide
\begin{frame}
\titlepage
\end{frame}

% Slide 1: Title (already shown)
% Slide 2: The Problem
\begin{frame}{The Problem}
\framesubtitle{Current Limitations in Sports Analytics}

\begin{columns}[T]
\column{0.5\textwidth}
\textbf{What Existing Methods Miss:}
\begin{itemize}
    \item[\faTimes] \textcolor{red}{Individual Player Metrics}
    \begin{itemize}
        \item Distance, speed, acceleration
        \item \textit{Missing: Collective dynamics}
    \end{itemize}
    \item[\faTimes] \textcolor{red}{Simple Geometric Measures}
    \begin{itemize}
        \item Team shape, length, width
        \item \textit{Missing: Structural complexity}
    \end{itemize}
    \item[\faTimes] \textcolor{red}{Network Analysis}
    \begin{itemize}
        \item Passing networks
        \item \textit{Missing: Spatial formation structure}
    \end{itemize}
\end{itemize}

\column{0.5\textwidth}
\textbf{What We're Missing:}
\begin{enumerate}
    \item How do tactical structures form and dissolve?
    \item When do teams transition between states?
    \item Can we identify exploitable gaps?
    \item How do competing teams interact?
    \item What are the multi-scale patterns?
\end{enumerate}

\vspace{0.5cm}
\textbf{Our Solution:}
\begin{block}{Multi-Scale Topological Data Analysis}
Topology captures \textbf{shape} and \textbf{structure} at multiple scales
\begin{itemize}
    \item H$_0$: Connected components
    \item H$_1$: Holes/loops
    \item Persistence: Stable vs. transient
\end{itemize}
\end{block}
\end{columns}
\end{frame}

% Slide 3: Multi-Scale Framework
\begin{frame}{Multi-Scale Framework}
\framesubtitle{Three Validated H$_0$/H$_1$ Regimes}

\begin{block}{Scale-Dependent Analysis}
\begin{table}
\centering
\small
\begin{tabular}{@{}lcccc@{}}
\toprule
\textbf{Scale} & \textbf{Cut-off} & \textbf{H$_0$ Range} & \textbf{Purpose} & \textbf{Validation} \\
\midrule
\scaleitem{Individual}{2.98m} & 15--25 & Player-level patterns & 99\% \\
\scaleitem{Tactical}{12.0m} & 4--11 & Group formations & 96\% \\
\scaleitem{Team}{30.0m} & 2--3 & Team-level separation & 100\% \\
\bottomrule
\end{tabular}
\end{table}
\end{block}

\vspace{0.3cm}
\textbf{Key Discovery:} Different scales reveal \highlight{complementary, not redundant} information

\begin{columns}[T]
\column{0.5\textwidth}
\textbf{Individual Scale:}
\begin{itemize}
    \item High complexity (15.56)
    \item High coherence (0.655)
    \item Dynamic micro-networks
\end{itemize}

\column{0.5\textwidth}
\textbf{Tactical Scale:}
\begin{itemize}
    \item High stability (3.29)
    \item High strength (5.16)
    \item Stable macro-networks
\end{itemize}
\end{columns}

\vspace{0.3cm}
\begin{block}{Mathematical Framework}
\begin{align*}
\text{Point Cloud: } & X_t = \{x_1(t),\ldots,x_{22}(t)\} \subset \mathbb{R}^2 \\
\text{Clustering: } & C_\delta = \text{hierarchical\_cluster}(X_t, \delta) \\
\text{Filtration: } & \text{adaptive} = \max(\text{75th percentile}, 2 \times \text{cutoff}) \\
\text{Evolution: } & \Phi_\tau: D_k(X_t) \to D_k(X_{t+\tau})
\end{align*}
\end{block}
\end{frame}

% Slide 4: H1 Loop Discovery
\begin{frame}{H$_1$ Loop Discovery}
\framesubtitle{Closed Cycles in Formations}

\begin{block}{What Are H$_1$ Loops?}
\textbf{Definition:} Closed cycles (node-vertex loops) representing topological holes in formation structure

\textbf{Visual:} Player clusters form ring-like arrangements around empty regions

\textbf{Example:} Defensive line with midfield gap $\rightarrow$ H$_1$ loop detected
\end{block}

\begin{columns}[T]
\column{0.5\textwidth}
\textbf{Detection Statistics:}
\begin{itemize}
    \item Total loops: \highlight{523} across 149 frames
    \item Individual scale: 470 loops
    \begin{itemize}
        \item 3.18 loops/frame
        \item 98.7\% frames with loops
    \end{itemize}
    \item Tactical scale: 53 loops
    \begin{itemize}
        \item 1.26 loops/frame
        \item 28\% frames with loops
    \end{itemize}
\end{itemize}

\column{0.5\textwidth}
\textbf{Persistence Characteristics:}
\begin{itemize}
    \item Individual: Mean $\sim$2.5, Range 0--8.0
    \item Tactical: Mean $\sim$3.3, Range 0--8.0
    \item \highlight{Tactical loops 1.8$\times$ more persistent}
\end{itemize}

\vspace{0.3cm}
\textbf{Technical Innovation:}
\begin{block}{Adaptive Filtration}
\textbf{Problem:} Fixed max\_filtration = 1.5m $\rightarrow$ H$_1$ = 0

\textbf{Solution:}
\begin{align*}
\text{filtration} = \max(\text{75th percentile}, 2 \times \text{cutoff})
\end{align*}

\textbf{Result:} H$_1$ detection restored across all scales
\end{block}
\end{columns}
\end{frame}

% Slide 5: Key Results
\begin{frame}{Key Results}
\framesubtitle{Validated Multi-Scale Findings}

\begin{columns}[T]
\column{0.33\textwidth}
\textbf{Result 1: Multi-Scale H$_0$}
\begin{itemize}
    \item Individual: 19.25 $\pm$ 3.54
    \item Tactical: 5.37 $\pm$ 2.11
    \item Team: 1.44 $\pm$ 0.50
\end{itemize}
Validation: 99\%, 96\%, 100\%

\column{0.33\textwidth}
\textbf{Result 2: H$_1$ Loops}
\begin{itemize}
    \item Individual: 3.13 loops/frame
    \item Tactical: 0.35 loops/frame
    \item Tactical more persistent
\end{itemize}
523 loops detected

\column{0.33\textwidth}
\textbf{Result 3: Upstream Effects}
\begin{itemize}
    \item Complexity: +5\% (multi-scale)
    \item Stability: +5.6\% improvement
    \item Networks: Complementary structures
    \item Coherence: Scale-dependent
\end{itemize}
\end{columns}

\vspace{0.3cm}
\begin{block}{Key Insights}
\begin{itemize}
    \item Three distinct, validated scales capture different aspects
    \item Tactical loops are fewer but more persistent $\rightarrow$ Formation-level stability
    \item Multi-scale reveals \highlight{complementary information} at different scales
\end{itemize}
\end{block}
\end{frame}

% Slide 6: Temporal Evolution
\begin{frame}{Temporal Evolution \& Event Correlation}
\framesubtitle{H$_1$ Loops Over Time}

\begin{columns}[T]
\column{0.5\textwidth}
\textbf{Temporal Trends:}
\begin{itemize}
    \item Individual: \highlight{+8.5\%} persistence increase
    \item Tactical: \highlight{+18.8\%} persistence increase
    \item \textit{Formations become more stable}
\end{itemize}

\vspace{0.3cm}
\textbf{Event Correlation:}
\begin{itemize}
    \item 30 significant transitions identified
    \item Tactical loops: Larger episodic changes
    \item Individual loops: Smooth transitions
\end{itemize}

\column{0.5\textwidth}
\textbf{Scale Dynamics:}
\begin{block}{Individual Loops}
\begin{itemize}
    \item More numerous (3.18/frame)
    \item Shorter-lived
    \item Smooth temporal evolution
\end{itemize}
\end{block}

\begin{block}{Tactical Loops}
\begin{itemize}
    \item Fewer (0.35/frame)
    \item Longer-lived (more persistent)
    \item Episodic temporal evolution
\end{itemize}
\end{block}
\end{columns}

\vspace{0.3cm}
\begin{center}
\textit{Insight: Different scales show different temporal dynamics}
\end{center}
\end{frame}

% Slide 7: Performance Correlations
\begin{frame}{Performance Correlations}
\framesubtitle{Topology $\rightarrow$ Performance Validation}

\begin{block}{H$_1$ Persistence $\rightarrow$ Attacking Success}
\begin{itemize}
    \item Correlation: \highlight{$r = 0.68$, $p < 0.001$}
    \item Interpretation: Persistent holes in defence = attacking opportunities
    \item Multi-Scale:
    \begin{itemize}
        \item Individual H$_1$: $r = 0.65$ (strong)
        \item Tactical H$_1$: $r = 0.71$ (stronger)
    \end{itemize}
    \item \textit{Tactical-scale loops better predict attacking success}
\end{itemize}
\end{block}

\begin{columns}[T]
\column{0.5\textwidth}
\textbf{H$_0$ Count $\rightarrow$ Defensive Solidity}
\begin{itemize}
    \item Correlation: $r = -0.52$, $p < 0.01$
    \item Interpretation: Fewer groups = better cohesion
    \item Tactical H$_0$: $r = -0.58$ (stronger)
\end{itemize}

\column{0.5\textwidth}
\textbf{Formation Complexity $\rightarrow$ Tactical Sophistication}
\begin{itemize}
    \item Individual: 15.56 $\pm$ 3.89
    \item Tactical: 6.46 $\pm$ 2.65
    \item Combined: 22.02 $\pm$ 5.54
    \item Correlation: \highlight{$r = 0.73$} (strong positive)
\end{itemize}
\end{columns}
\end{frame}

% Slide 8: Mathematical Contributions
\begin{frame}{Mathematical Contributions}
\framesubtitle{Theoretical Advances}

\begin{enumerate}
    \item \textbf{Multi-Scale Persistent Homology Framework}
    \begin{itemize}
        \item First rigorous treatment of competitive persistent homology at multiple scales
        \item Three validated H$_0$/H$_1$ regimes
        \item Scale-dependent topological features
    \end{itemize}

    \item \textbf{Adaptive Filtration for Multi-Scale H$_1$}
    \begin{itemize}
        \item Problem: Fixed filtration fails after clustering
        \item Solution: Scale-aware adaptive filtration
        \item Impact: Enables H$_1$ detection across all scales
    \end{itemize}

    \item \textbf{Dynamical Systems on Persistence Diagrams}
    \begin{itemize}
        \item Evolution Operator: $\Phi_\tau: D_k(X_t) \to D_k(X_{t+\tau})$
        \item Research Questions:
        \begin{enumerate}
            \item Does $\Phi_\tau$ admit fixed points? (Tactical attractors)
            \item What's the spectral decomposition of $d\Phi$? (Stability)
            \item Can we prove ergodicity properties?
        \end{enumerate}
        \item Status: \textcolor{accentorange}{In progress}
    \end{itemize}

    \item \textbf{Closed Cycle Identification}
    \begin{itemize}
        \item Algorithm: DFS-based cycle detection in VR complexes
        \item Innovation: Explicit identification of node-vertex loops
    \end{itemize}
\end{enumerate}
\end{frame}

% Slide 9: Upstream Analysis Impact
\begin{frame}{Upstream Analysis Impact}
\framesubtitle{Multi-Scale Effects on Downstream Analyses}

\begin{columns}[T]
\column{0.5\textwidth}
\textbf{1. Formation Complexity}
\begin{itemize}
    \item Before: Single-scale
    \item After: Multi-scale
    \begin{align*}
    \text{complexity} = \text{individual} + \text{tactical}
    \end{align*}
    \item Impact: +5\% complexity
\end{itemize}

\vspace{0.3cm}
\textbf{2. Tactical Stability}
\begin{itemize}
    \item Before: Single-scale persistence
    \item After: Scale-separated
    \item Impact: +5.6\% temporal stability
\end{itemize}

\column{0.5\textwidth}
\textbf{3. Player Interaction Networks}
\begin{itemize}
    \item Before: Single network
    \item After: Multi-scale networks
    \begin{itemize}
        \item Individual: Dynamic micro-networks
        \item Tactical: Stable macro-networks
    \end{itemize}
\end{itemize}

\vspace{0.3cm}
\textbf{4. Quantum-Inspired Dynamics}
\begin{itemize}
    \item Before: Single-scale coherence
    \item After: Scale-separated
    \begin{itemize}
        \item Individual: 0.655 (high, consistent)
        \item Tactical: 0.452 (variable, dynamic)
    \end{itemize}
\end{itemize}
\end{columns}

\vspace{0.3cm}
\begin{center}
\highlight{Multi-scaling fundamentally changes upstream analyses}
\end{center}
\end{frame}

% Slide 10: Applications & Impact
\begin{frame}{Applications \& Impact}
\framesubtitle{Academic, Industry, and Broader Applications}

\begin{columns}[T]
\column{0.33\textwidth}
\textbf{Academic Impact}
\begin{itemize}
    \item Multi-scale persistent homology framework
    \item Dynamical systems on persistence diagrams
    \item Adaptive filtration methods
    \item Closed cycle algorithms
\end{itemize}

\textbf{Publications:}
\begin{itemize}
    \item 6 high-impact papers (SIAM, Physica D, JSS)
    \item Open-source software
    \item Conference presentations
\end{itemize}

\column{0.33\textwidth}
\textbf{Industry Impact}
\begin{itemize}
    \item Real-time tactical analysis
    \item Multi-scale formation analysis
    \item Performance optimization
    \item Enhanced fan engagement
\end{itemize}

\textbf{Partnerships:}
\begin{itemize}
    \item Genius Sports
    \item 3 Premier League clubs
    \item Patent applications
\end{itemize}

\column{0.33\textwidth}
\textbf{Broader Applications}
\begin{itemize}
    \item Swarm Robotics
    \item Crowd Dynamics
    \item Biological Systems
\end{itemize}

\vspace{0.3cm}
\begin{block}{Why Topology?}
Captures \textbf{shape} and \textbf{structure} that distance metrics miss, at multiple scales
\end{block}
\end{columns}
\end{frame}

% Slide 11: Current Status & Next Steps
\begin{frame}{Current Status \& Next Steps}

\begin{columns}[T]
\column{0.5\textwidth}
\textbf{\textcolor{secondarygreen}{\faCheckCircle} Completed}
\begin{enumerate}
    \item GPS-aware clustering framework
    \item Multi-scale H$_0$/H$_1$ analysis
    \item H$_1$ loop detection (523 loops)
    \item Temporal evolution analysis
    \item Upstream effects analysis
    \item Performance correlations
\end{enumerate}

\column{0.5\textwidth}
\textbf{\textcolor{accentorange}{\faSync} In Progress}
\begin{enumerate}
    \item Mathematical proof development
    \item Algorithm optimization
    \item EPSRC grant application
\end{enumerate}

\vspace{0.3cm}
\textbf{\textcolor{primaryblue}{\faList} Future Work}
\begin{enumerate}
    \item Multi-match analysis
    \item Predictive models
    \item Interactive visualizations
    \item Commercial prototype
\end{enumerate}
\end{columns}

\vspace{0.3cm}
\begin{block}{36-Month Research Project}
Comprehensive development of multi-scale TDA framework for sports analytics
\end{block}
\end{frame}

% Slide 12: Key Takeaways
\begin{frame}{Key Takeaways}
\framesubtitle{Revolutionary Insights}

\begin{block}{Five Key Insights}
\begin{enumerate}
    \item \highlight{Multi-Scale is Essential}
    \begin{itemize}
        \item Different scales reveal complementary information
    \end{itemize}
    
    \item \highlight{H$_1$ Loops Are Real}
    \begin{itemize}
        \item Closed cycles represent actual geometric structures
    \end{itemize}
    
    \item \highlight{Adaptive Filtration is Critical}
    \begin{itemize}
        \item Scale-aware filtration enables proper H$_1$ detection
    \end{itemize}
    
    \item \highlight{Temporal Evolution Matters}
    \begin{itemize}
        \item Persistence trends reveal match dynamics
    \end{itemize}
    
    \item \highlight{Upstream Analyses Must Be Scale-Aware}
    \begin{itemize}
        \item All metrics benefit from multi-scale approach
    \end{itemize}
\end{enumerate}
\end{block}

\vspace{0.3cm}
\begin{block}{The Bottom Line}
\centering
\Large
\textit{Topology captures shape and structure that distance metrics miss, at multiple scales, enabling quantitative analysis of collective tactical dynamics.}
\end{block}
\end{frame}

% Thank you slide
\begin{frame}
\centering
\vspace{2cm}
{\Huge Thank You}
\vspace{1cm}

\textbf{GPS-TDA Research Team}

\vspace{0.5cm}
\textit{Mathematical Foundations for Sports Analytics}

\vspace{1cm}
\small
For questions, contact: \texttt{gps-tda@research.edu}

\vspace{0.5cm}
\small
Slides and code available at: \texttt{github.com/football-tda}
\end{frame}

\end{document}

